\begin{table}[tp]
\centering
\caption{Necessary oracle budgets for \href{https://github.com/k3nfalt/longtails/blob/main/LongTails/Setups/GeneralizedSmoothConvexPolynomialDimension.lean}{GeneralizedSmoothConvexPolynomialDimension}.}\label{tab:GeneralizedSmoothConvexPolynomialDimension}
\parbox{\linewidth}{\footnotesize Convex twice differentiable $f:\mathbb{R}^d\to\mathbb{R}$ satisfies $\|\nabla^2 f(x)\|\le L_0+L_1\|\nabla f(x)\|$ and has a minimizer $x^\star$ with $\|x_0-x^\star\|\le R$; no initial-gap bound. Parameters are positive, $q=L_1R\ge 2^{20}$, $L_0/L_1^2\le\varepsilon\le L_0R^2/16$, and $d\ge\lceil 2^{44}(1+q)^5\rceil$. Algorithms may be arbitrary randomized adaptive exact-gradient algorithms. The target is an integrable objective gap with $\int [f(\hat x(\omega))-f(x^\star)]\,dP(\omega)\le\varepsilon$, where $P$ is the algorithm's probability law.}\par\smallskip
\begin{threeparttable}
\footnotesize
\setlength{\tabcolsep}{4pt}\renewcommand{\arraystretch}{1.25}
\begin{tabularx}{\linewidth}{@{}>{\raggedright\arraybackslash}p{0.20\linewidth}>{\raggedright\arraybackslash}p{0.32\linewidth}>{\raggedright\arraybackslash}p{0.25\linewidth}>{\raggedright\arraybackslash}X@{}}
\toprule
\textbf{Result} & \textbf{Oracle complexity} & \textbf{Novelty} & \textbf{Limitations} \\
\midrule
\textbf{(1)} \href{https://github.com/k3nfalt/longtails/blob/main/LongTails/Theorems/PolynomialDimensionGeneralizedSmoothLowerBound.lean}{Polynomial\allowbreak Dimension\allowbreak Generalized\allowbreak Smooth\allowbreak Lower\allowbreak Bound} \textbf{(Ours)} $(\star\star)$ \href{https://github.com/k3nfalt/longtails/blob/main/scripts/reviews/PolynomialDimensionGeneralizedSmoothLowerBound.md}{(review)}
& Lower bound (necessary): $N=\Omega(q/\ell)$ exact-gradient calls; fixed-budget failure on a seed-independent instance. & & \\
\midrule[1.2pt]
\textbf{(2)} \citep{mirrorEverestV2}, Thms.~4.8, 4.10
& Lower bound (necessary): $N=\Omega\!\left(1+\sqrt{L_0R^2/\varepsilon}+q/\sqrt{\ell}\right)$ value-and-gradient calls. Deterministic algorithms; $\varepsilon\le L_0R^2$.
& \textbf{(1)}: Scope: hardness extends from deterministic to arbitrary randomized algorithms, with expected-gap accuracy.
& \textbf{(1)}: Weaker $q$ term by $\sqrt{\ell}$; no separate $\sqrt{L_0R^2/\varepsilon}$ term. Gradient-only hardness; $d\ge\lceil2^{44}(1+q)^5\rceil$ versus $\mathcal{O}(q)$ in the shared regime. Requires $q\ge2^{20}$ and $\varepsilon\ge L_0/L_1^2$. \\
\midrule
\textbf{(3)} \citep{polynomialPriorRepository166}, repository result, commit 0ed1e4a6
& Lower bound (necessary): $N=\Omega(q/\ell)$ exact-gradient calls. Randomized, fixed-budget expected-gap hardness; $d\ge\lceil2^{20}(1+q)^8e^{2\sqrt q}\rceil$; same parameter regime.
& \textbf{(1)}: Scope: required dimension decreases from $\lceil2^{20}(1+q)^8e^{2\sqrt q}\rceil$ to $\lceil2^{44}(1+q)^5\rceil$. Matches the necessary budget and randomized expected-gap criterion.
& \\
\midrule
\textbf{(4)} \citep{polynomialEverestIIV2}, Thm.~7.1; different finite-sum class
& Lower bound (necessary expected calls):
$\begin{aligned}
\Omega\!\Bigl(\bigl[n(1+q^{2/3})&{}\\
{}+\sqrt{n(L_0+L_1^2\delta)R^2/\varepsilon+nq^2}&\bigr]/\Lambda^2\Bigr).
\end{aligned}$
Component value-and-gradient oracle; convex bounded-below $C^2$ components, $\|\nabla^2f_i\|\le L_{0,i}+L_1\|\nabla f_i\|$, $n^{-1}\sum_i L_{0,i}\le L_0$; $n\ge2$, arbitrary dimension, $\Delta\ge2\varepsilon$, $q^2\delta/\varepsilon\ge3$.
& \textbf{(1)}: Single-objective hardness in polynomial dimension. Finite-sum classes are incomparable: their average need not satisfy the single-function curvature bound. Rates are not comparable; $n=1$ is excluded.
& \textbf{(1)}: Gradient-only, fixed-budget failure; \textbf{(4)} permits component values and gradients and bounds expected stopping cost. \\
\bottomrule
\end{tabularx}
\begin{tablenotes}[flushleft]\footnotesize
\item[] \textbf{Abbr:} $N$: necessary calls; $q=L_1R$; $\ell=1+\max\{0,\log q\}$. For \textbf{(4)}, $f=n^{-1}\sum_{i=1}^n f_i$ denotes the finite-sum objective; $\Delta$ bounds $f(x_0)-f^\star$, $\delta$ bounds $f^\star-n^{-1}\sum_i\inf f_i$, and $\Lambda=1+\log\!\left(n+(L_0+L_1^2\delta)R^2/\varepsilon+q^2\right)$. $\Omega$ denotes a necessary budget up to universal factors; larger lower bounds are stronger.
\end{tablenotes}
\end{threeparttable}
\end{table}
\clearpage
